\subsection{\unclass Assumptions, Limitations, and Gaps}  

\guidance{
	\unclass TIMELINE: This section should be started at \gls{srr} and finalized by \gls{cdr}. Please note that this is expected to be a collaborative effort between the intel analyst, \gls{tpoc}, and model leads with back forth.  
	
	\unclass Threat assumptions are necessary to complete model development when certain information is needed but insufficient intelligence data exists.  Such assumptions are often labeled “engineering judgment.”  Documentation should identify the basis for key assumptions. The intelligence products we have access to are themselves derived from a number of sub-sources and assumptions.   Documentation will identify key assumptions for the defined use cases evaluated in the \gls{vv} report.  Documenting all assumptions will enable users to perform a key assumptions check for use cases other than the defined use cases identified in the \gls{vv} report.
	
	\unclass It is essential that the assumptions used in the model are properly captured and documented to allow the user to understand the reasons to use (or not use) alternative assumptions in the future, and to understand where the model is more limited in emulating the associated system.  It is also imperative for future users when using the model in ways that were not tested or documented during model development, as future testing will need to observe and understand the model's assumptions in order for the model to continue to behave properly outside of the \gls{vv} report's scope.
	
	\unclass Intelligence policy requires clearly distinguishing assumptions and judgments from underlying intelligence information. Documentation will comply with distinctions and articulate implications if key assumptions are incorrect.  The list of assumptions will enable users to identify linchpin assumptions for uses outside of the evaluated test case(s).
	
	\unclass More details: 
	
	\unclass Discuss assumptions, limitations, and gaps including: gaps in intelligence, limited implementations of intelligence, assumptions that either fill in the gaps above or simplify model behavior compared to intelligence information, gaps in implementation due to time/funding prioritization, assumptions on approximate code solutions to complex intelligence issues, and limitations in the code base for the model that lead to changes in behavior.  
	
	\unclass List and discuss how intelligence gaps affect model development, how the gaps were filled and what assumptions were used, and the potential impact or limitations to model usefulness. Specific component details should have these captured in their respective sections.  

}